%-------------------------
% Logan MacAskill - resume
% Jake's Resume template (MIT), https://github.com/jakegut/resume
% Ported from portfolio/public/docs/resume.md on 2026-09-12
%------------------------

% FONT SIZE. Plain \documentclass{article} only accepts 10, 11 and 12pt, so
% getting under 10 means the extarticle class, which adds 8pt and 9pt. Overleaf
% has extsizes installed, nothing to do.
%
% 9pt is about the floor for a resume. Below that it reads as someone who could
% not bear to cut anything. If it fits at 9.5 or 10, use that instead.
\documentclass[letterpaper,9pt]{extarticle}
% \documentclass[letterpaper,10pt]{article}   % the original, swap back freely

\usepackage{latexsym}
\usepackage[empty]{fullpage}
\usepackage{titlesec}
\usepackage{marvosym}
% xcolor rather than color, so the dvipsnames palette is available for links.
% Drop-in replacement, nothing else changes.
\usepackage[usenames,dvipsnames]{xcolor}
\usepackage{verbatim}
\usepackage{enumitem}

% LINKS. Underline alone reads as emphasis, not as something clickable, so
% every link is also colored. MidnightBlue is dark enough to stay readable on
% paper and in grayscale; a saturated blue looks like a 1998 web page.
%
% The underline stays on purpose. Color dies when someone prints in grayscale,
% the underline survives, so the two together cover both cases.
%
% Color has no effect on ATS parsing. Text extraction ignores it entirely.
\usepackage[colorlinks=true, urlcolor=MidnightBlue, linkcolor=MidnightBlue]{hyperref}
% \usepackage[hidelinks]{hyperref}   % the original, all-black links
\usepackage{fancyhdr}
\usepackage[english]{babel}
\usepackage{tabularx}
\input{glyphtounicode}

\pagestyle{fancy}
\fancyhf{}
\fancyfoot{}
\renewcommand{\headrulewidth}{0pt}
\renewcommand{\footrulewidth}{0pt}

% Margins. Widened from Jake's defaults (which give 0.5in sides) to 0.4in.
% This is aimed at your actual complaint: a bullet that spills three words onto
% a second line costs a whole line, so widening the text block kills those
% specific wraps without shrinking any type. 0.4in is the floor. Below it some
% print paths clip the edge.
\addtolength{\oddsidemargin}{-0.6in}
\addtolength{\evensidemargin}{-0.6in}
\addtolength{\textwidth}{1.2in}
\addtolength{\topmargin}{-.6in}
\addtolength{\textheight}{1.2in}

\urlstyle{same}

\raggedbottom
\raggedright
\setlength{\tabcolsep}{0in}

% Sections formatting
\titleformat{\section}{
  \vspace{-4pt}\scshape\raggedright\large
}{}{0em}{}[\color{black}\titlerule \vspace{-5pt}]

% This line is the whole reason for the switch. It embeds a ToUnicode map, so
% every glyph in the PDF maps back to a real character. That is what the
% autofill parsers read.
\pdfgentounicode=1

% SPACING. Two separate knobs, deliberately pulling opposite directions.
%
% \linespread tightens the lines INSIDE one bullet, so a bullet that wraps
% reads as a single block instead of two loose lines. 0.95 is safe at 9pt;
% below about 0.90 descenders start crowding the line beneath.
%
% itemsep (down in \resumeItemListStart) opens the gap BETWEEN bullets. That
% is what makes them separable at a glance once the inner lines are tight.
%
% Net cost is a few lines: the gaps add more than the tightening saves. Drop
% itemsep to 2pt first if it pushes you over.
\linespread{0.95}

%-------------------------
% Custom commands

% The \vspace{-2pt} that used to live INSIDE this macro is the blank-line bug.
% A \vspace inside the item's own paragraph has to attach to something, so
% when a bullet's last line happens to fill the measure exactly, TeX opens a
% fresh line to hold the glue and you get a phantom empty line under it. The
% spacing now lives on the list instead, where it cannot do that.
\newcommand{\resumeItem}[1]{
  \item\small{#1}
}

\newcommand{\resumeSubheading}[4]{
  \vspace{-2pt}\item
    \begin{tabular*}{0.97\textwidth}[t]{l@{\extracolsep{\fill}}r}
      \textbf{#1} & #2 \\
      \textit{\small#3} & \textit{\small #4} \\
    \end{tabular*}\vspace{-7pt}
}

\newcommand{\resumeSubSubheading}[2]{
    \item
    \begin{tabular*}{0.97\textwidth}{l@{\extracolsep{\fill}}r}
      \textit{\small#1} & \textit{\small #2} \\
    \end{tabular*}\vspace{-7pt}
}

\newcommand{\resumeProjectHeading}[2]{
    \item
    \begin{tabular*}{0.97\textwidth}{l@{\extracolsep{\fill}}r}
      \small#1 & #2 \\
    \end{tabular*}\vspace{-7pt}
}

\newcommand{\resumeSubItem}[1]{\resumeItem{#1}\vspace{-4pt}}

\renewcommand\labelitemii{$\vcenter{\hbox{\tiny$\bullet$}}$}

\newcommand{\resumeSubHeadingListStart}{\begin{itemize}[leftmargin=0.15in, label={}]}
\newcommand{\resumeSubHeadingListEnd}{\end{itemize}}
% itemsep is the gap between bullets. topsep is the gap above the first one.
% parsep and partopsep are set to zero so the only vertical space in this list
% is the one you are actually controlling.
\newcommand{\resumeItemListStart}{\begin{itemize}[leftmargin=0.2in, itemsep=3pt, parsep=0pt, topsep=3pt, partopsep=0pt]}
\newcommand{\resumeItemListEnd}{\end{itemize}\vspace{-3pt}}

%-------------------------------------------
%%%%%%  RESUME STARTS HERE  %%%%%%%%%%%%%%%%%%%%%%%%%%%%

\begin{document}

%----------HEADING----------
\begin{center}
    \textbf{\Huge \scshape Logan MacAskill} \\ \vspace{3pt}
    \small Los Altos, CA $|$ (650) 237-9593 $|$
    \href{mailto:logan@macaskill.com}{\underline{logan@macaskill.com}} \\ \vspace{2pt}
    \href{https://logan.macaskill.com}{\underline{logan.macaskill.com}} $|$
    \href{https://github.com/log2bits}{\underline{github.com/log2bits}} $|$
    \href{https://www.linkedin.com/in/logan-macaskill-356190222/}{\underline{linkedin.com/in/logan-macaskill}}
\end{center}

%-----------ABOUT ME-----------
\section{About Me}
 \begin{itemize}[leftmargin=0.15in, label={}]
    \small{\item{Third-year CS student with two software engineering internships spent shipping and tuning production systems. I love problem solving, especially when I can take complete advantage of the hardware. I have worked across real-time systems, graphics, AI, and data science, and I am currently building my own Vulkan and C++ deferred renderer.}}
 \end{itemize}

%-----------EDUCATION-----------
\section{Education}
  \resumeSubHeadingListStart
    \resumeSubheading
      {University of California, Santa Cruz}{Santa Cruz, CA}
      {B.S. in Computer Science, GPA 3.6/4.0}{Expected Jun 2028}
  \resumeSubHeadingListEnd

%-----------EXPERIENCE-----------
\section{Experience}
  \resumeSubHeadingListStart

    \resumeSubheading
      {Software Engineering Intern}{Jun 2023 -- Aug 2023}
      {Flickr}{Mountain View, CA}
      \resumeItemListStart
        \resumeItem{Cut compute cost 92\% and raised throughput 39\% on a Python AWS Lambda service repairing 10 billion photos.}
        \resumeItem{Tuned memory and thread allocation to the cheapest stable configuration.}
        \resumeItem{Profiled dozens of Lambda configurations in Splunk, modeling cost per job, throughput, and tail latency.}
        \resumeItem{Recovered AWS's undocumented vCPU-to-memory scaling ratio by linear regression.}
      \resumeItemListEnd

    \resumeSubheading
      {Software Engineering Intern}{Jun 2022 -- Aug 2022}
      {SmugMug}{Mountain View, CA}
      \resumeItemListStart
        \resumeItem{Owned a user-facing gallery stats feature end to end in PHP, React and Next.js, from PR review through QA.}
        \resumeItem{Added daily refresh scheduling, clearer labels and empty states to the stats UI.}
        \resumeItem{Wrote unit tests and backend changes for a release-critical PHP 8.1 upgrade, landing it on schedule.}
      \resumeItemListEnd

  \resumeSubHeadingListEnd

%-----------PROJECTS-----------
\section{Projects}
    \resumeSubHeadingListStart

      \resumeProjectHeading
          {\href{https://logan.macaskill.com/experience/ray-vox}{\textbf{\underline{ray-vox: Ray-traced Voxel Renderer}}} $|$ \emph{Rust, WebGPU, Data Structures, Optimization}}{2026}
          \resumeItemListStart
            \resumeItem{Packed an 84 MB voxel scene into 16 MB with a custom sparse-voxel structure, beating zstd at max (19 MB).}
            \resumeItem{Rays walk the compressed form directly, with no decompression step.}
            \resumeItem{Removed every child pointer with per-node bitmasks and popcount, packing each node offset into one 32-bit word.}
            \resumeItem{Traversal stays branch-light, so GPU warps stay coherent.}
            \resumeItem{Designed it as a GPU acceleration structure with level-of-detail and matching on-disk and in-GPU layouts.}
          \resumeItemListEnd

      \resumeProjectHeading
          {\href{https://saisgonerogue.itch.io/sanity-check}{\textbf{\underline{Sanity Check: Logic Puzzle Game}}} $|$ \emph{Unity, C\#, Constraint Solving, Procedural Generation}}{2026}
          \resumeItemListStart
            \resumeItem{Built the puzzle engine for a five-person, one-week game jam.}
            \resumeItem{Encoded every possible answer as one bit in a hand-written bitset, so solving a room is set intersection.}
            \resumeItem{Generated puzzles that are provably unique and minimal, rejecting any room where a clue could be dropped.}
            \resumeItem{Hand-wrote the set layer with 64-bit words and SWAR popcount, scoring 20,000 candidate rooms per puzzle.}
            \resumeItem{Tuned five difficulty tiers on how many statements must combine before the first deduction.}
            \resumeItem{Also owned the game's lighting, post-processing and camera feel.}
          \resumeItemListEnd

      \resumeProjectHeading
          {\href{https://logan.macaskill.com/experience/spectral-raymarcher}{\textbf{\underline{Real-time Dielectric Spectral Raymarcher}}} $|$ \emph{GLSL, Spectral Rendering, Sampling, GPU}}{2026}
          \resumeItemListStart
            \resumeItem{Rendered a physically based dispersive diamond in a single real-time GLSL shader, live in a browser tab.}
            \resumeItem{Resolved 65,536 wavelengths from two samples per pixel by scheduling them across a 16x16 Bayer dither.}
            \resumeItem{Replaced raymarching with exact ray-plane intersection across the icosahedron's 20 faces.}
            \resumeItem{Peeled exact Fresnel energy per facet instead of sampling one path, staying noise-free at one sample.}
            \resumeItem{Terminated each ray below 4\% trapped energy.}
            \resumeItem{Kept every pixel a self-contained shader invocation, with no frame history.}
          \resumeItemListEnd

      \resumeProjectHeading
          {\href{https://logan.macaskill.com/experience/crowd-surfers}{\textbf{\underline{Crowd Surfers: Real-time 3D Game}}} $|$ \emph{Godot, Shaders, Lighting, Architecture}}{2025 -- 2026}
          \resumeItemListStart
            \resumeItem{Convinced a 100-person student team to rebuild the renderer mid-project by building a prototype instead of arguing.}
            \resumeItem{Moved it from faked-depth sprites to true 3D with an angled orthographic camera.}
            \resumeItem{Built an occlusion shader fading buildings as the player skates behind them, with a frustum test and dithered alpha.}
          \resumeItemListEnd

      \resumeProjectHeading
          {\href{https://logan.macaskill.com/experience/coaxial-swerve-drive}{\textbf{\underline{Coaxial Swerve Drive}}} $|$ \emph{Java, Control Theory, Computer Vision}}{2023 -- 2024}
          \resumeItemListStart
            \resumeItem{Led an off-season FRC team to a working coaxial swerve prototype in one month, four months ahead of schedule.}
            \resumeItem{Built the CNC chassis, the electronics, and 2,000 lines of Java.}
            \resumeItem{Held sub-centimeter localization through wheel slip by fusing encoders, gyro and AprilTag vision in a Kalman filter.}
            \resumeItem{Ran per-module PID and feedforward on a 1 kHz control loop after self-teaching graduate-level control theory.}
            \resumeItem{Simulated torque curves, inertia and brownouts on the production code, deriving current limits from tire friction.}
          \resumeItemListEnd

    \resumeSubHeadingListEnd

%-----------TECHNICAL SKILLS-----------
\section{Technical Skills}
 \begin{itemize}[leftmargin=0.15in, label={}]
    \small{\item{
     \textbf{Languages}{: Rust, C++, C\#, Python, Java, TypeScript / JavaScript, PHP, GLSL / WGSL} \\
     \textbf{Graphics \& GPU}{: Vulkan, WebGPU, Ray Tracing, Real-Time Rendering, Rasterization, Shaders, Level-of-Detail} \\
     \textbf{Tools \& Cloud}{: Git, Linux, AWS (EC2, Lambda, Graviton), Docker, OpenCV, OpenAI API, Godot, Unity, React / Next.js} \\
     \textbf{Concepts}{: Performance Optimization, GPU Compute, Computer Vision, Pose Estimation, LLMs / Prompt Engineering, Control Theory}
    }}
 \end{itemize}

%-------------------------------------------
\end{document}
